


\begin{lstlisting}[language=Python]
import numpy as np
import pandas as pd
from sklearn.model_selection import train_test_split, KFold
from sklearn.linear_model import Ridge, LinearRegression
from sklearn.kernel_ridge import KernelRidge
from sklearn.pipeline import make_pipeline
from sklearn.preprocessing import StandardScaler
from tqdm import tqdm
from sklearn.metrics import mean_squared_error

number_of_folds = 10
data = pd.read_csv("../data/bodyfat.csv")
X = data.drop(columns=["BodyFat", "Density"])
y = data["BodyFat"]
X_train, X_test, y_train, y_test = train_test_split(
    X, y, test_size=0.2, random_state=0
    )
\end{lstlisting}
\newpage
\subsubsection*{(i) Erro médio quadrático da regressão linear}

\begin{lstlisting}[language=Python]
linear_regression_model = LinearRegression(fit_intercept=True)
linear_regression_model.fit(X_train, y_train)
y_hat = linear_regression_model.predict(X_test)
mse_linear = mean_squared_error(y_true=y_test, y_pred=y_hat)

y_hat_train = linear_regression_model.predict(X_train)
mse_linear_train = mean_squared_error(y_true=y_train, y_pred=y_hat_train)
\end{lstlisting}
\begin{table}[h!]
    \centering
    \begin{tabular}{|l|c|}
        \hline
        Conjunto & Erro médio quadrático       \\ \hline
        Teste    & \textbf{15.910489104854445} \\
        Treino   & \textbf{18.090133641655044} \\
        \hline
    \end{tabular}
    \caption{Resultados da regressão linear}
\end{table}
\newpage
\subsubsection*{(ii) Erro médio quadrático da regressão ridge}

\begin{lstlisting}[language=Python]
alphas = 10 ** np.linspace(3, -2, 100)
ridge_cv_mse = []
for alpha in tqdm(alphas):
    cv_MSE = []
    folds = KFold(
        n_splits=number_of_folds, shuffle=True, random_state=2023
    ).split(X_train, y_train)

    for train_idx, val_idx in folds:
        ridge_pipeline = make_pipeline(
            StandardScaler(), Ridge(alpha=alpha)
        )
        ridge_pipeline.fit(
            X_train.iloc[train_idx], y_train.iloc[train_idx]
        )
        y_hat = ridge_pipeline.predict(
            X_train.iloc[val_idx]
        )
        cv_MSE.append(
            mean_squared_error(y_true=y_train.iloc[val_idx], y_pred=y_hat)
        )

    ridge_cv_mse.append(np.mean(cv_MSE))

optimal_alpha = alphas[np.argmin(ridge_cv_mse)]
ridge_pipeline = make_pipeline(
    StandardScaler(), Ridge(alpha=optimal_alpha)
)
ridge_pipeline.fit(X_train, y_train)
y_hat_ridge = ridge_pipeline.predict(X_test)
ridge_test_mse = mean_squared_error(y_true=y_test, y_pred=y_hat_ridge)

y_hat_ridge_train = ridge_pipeline.predict(X_train)
ridge_train_mse = mean_squared_error(
    y_true=y_train, y_pred=y_hat_ridge_train
)
\end{lstlisting}

O melhor valor de $\alpha$ encontrado foi \textbf{0.52}. Com ele, obteve-se os seguintes erros médios
quadráticos:

\begin{table}[h!]
    \centering
    \begin{tabular}{|l|c|}
        \hline
        Conjunto             & Erro médio quadrático \\ \hline
        Validação (10 folds) & \textbf{21.84}        \\
        Teste                & \textbf{15.87}        \\
        Treino               & \textbf{18.099}       \\
        \hline
    \end{tabular}
    \caption{Resultados da regressão ridge}
\end{table}
\newpage
\subsubsection*{(iii) Kernel para Ridge Regression}

\begin{lstlisting}[language=Python]
kernels = ["linear", "polynomial", "rbf", "laplacian"]
gammas = [10**-3]
alphas = 10 ** np.linspace(3, -2, 100)
hyperparams = [
    (kernel, gamma, alpha)
    for kernel in kernels
    for gamma in gammas
    for alpha in alphas
]

kernel_ridge_cv_mse = []
for hyperparam in tqdm(hyperparams):
    kernel_fold, gamma_fold, alpha_fold = hyperparam
    cv_MSE = []
    folds = KFold(
        n_splits=number_of_folds, shuffle=True, random_state=2023
    ).split(X_train, y_train)

    for train_idx, val_idx in folds:
        kernel_ridge_pipeline = make_pipeline(
            StandardScaler(),
            KernelRidge(
                kernel=kernel_fold, 
                alpha=alpha_fold,
                gamma=gamma_fold
            ),
        )
        kernel_ridge_pipeline.fit(
            X_train.iloc[train_idx], y_train.iloc[train_idx]
        )
        y_hat = kernel_ridge_pipeline.predict(X_train.iloc[val_idx])
        cv_MSE.append(
            mean_squared_error(y_true=y_train.iloc[val_idx], y_pred=y_hat)
        )

    kernel_ridge_cv_mse.append(np.mean(cv_MSE))

optimal_hyperparam = hyperparams[np.argmin(kernel_ridge_cv_mse)]
optimal_kernel, optimal_gamma, optimal_alpha = optimal_hyperparam

kernel_ridge_pipeline = make_pipeline(
    StandardScaler(),
    KernelRidge(
        kernel=optimal_kernel, alpha=optimal_alpha, gamma=optimal_gamma
    ),
)
kernel_ridge_pipeline.fit(X_train, y_train)
kernel_y_hat_ridge = kernel_ridge_pipeline.predict(X_test)
kernel_ridge_test_mse = mean_squared_error(
    y_true=y_test, y_pred=kernel_y_hat_ridge
)

kernel_y_hat_ridge_train = kernel_ridge_pipeline.predict(X_train)
kernel_ridge_train_mse = mean_squared_error(
    y_true=y_train, y_pred=kernel_y_hat_ridge_train
)
\end{lstlisting}

O melhor conjunto de hiperparâmetros encontrado é:
este é: 14.66379896222263
reino é: 17.223450675066484

\begin{table}[h!]
    \centering
    \begin{tabular}{|c|c|}
        \hline
        Hiperparâmetro & Valor          \\ \hline
        $\alpha$       & \textbf{0.01}  \\
        $\gamma$       & \textbf{0.001} \\
        kernel         & \textbf{rbf}   \\ \hline
    \end{tabular}
\end{table}

Com ele, os erros médios encontrados foram:
\begin{table}[h!]
    \centering
    \begin{tabular}{|l|c|}
        \hline
        Conjunto             & Erro médio quadrático \\ \hline
        Validação (10 folds) & \textbf{20.55}        \\
        Teste                & \textbf{14.66}        \\
        Treino               & \textbf{17.22}        \\
        \hline
    \end{tabular}
    \caption{Resultados da regressão ridge com kernel}
\end{table}


\newpage
\subsubsection*{(iv) Comparação dos modelos}

\textbf{Resultados no conjunto de treino:}
\begin{table}[h!]
    \centering
    \begin{tabular}{|l|c|}
        \hline
        Modelo                     & Erro médio quadrático \\ \hline
        Regressão linear           & \textbf{18.09}        \\
        Regressão ridge            & \textbf{18.53}        \\
        Regressão ridge com kernel & \textbf{17.22}        \\
        \hline
    \end{tabular}
    \caption{Resultados de treino}
\end{table}

\textbf{Resultados no conjunto de teste:}
\begin{table}[h!]
    \centering
    \begin{tabular}{|l|c|}
        \hline
        Modelo                     & Erro médio quadrático \\ \hline
        Regressão linear           & \textbf{15.91}        \\
        Regressão ridge            & \textbf{15.87}        \\
        Regressão ridge com kernel & \textbf{14.66}        \\
        \hline
    \end{tabular}
    \caption{Resultados de teste}
\end{table}

Dos três modelos analizados, o que apresentou melhor desempenho foi a regressão ridge com kernel, tanto no conjunto de treino quanto no conjunto de teste. A
regressão linear e a regressão ridge apresentaram desempenhos semelhantes, com a regressão ridge tendo um desempenho ligeiramente pior no conjunto de treino,
mas melhor no conjunto de teste. Isso sugere que a regressão ridge pode ter ajudado a reduzir o overfitting em comparação com a regressão linear,
enquanto a adição do kernel na regressão ridge permitiu capturar relações não lineares nos dados, resultando em um desempenho superior.